\section{Conclusion}
\label{sec:conclusion}
This work evaluated static prediction of flaky tests in C\# from code vocabulary. The lexical signal exists, but it is modest. Across the 17 projects, which total 19,211 executed test cases, corresponding to 12,665 labeled test methods, of which 30 are flaky, the best models reach an MCC of 0.51--0.53 and an AUC of 0.86--0.87 under cross-validation, only moderately above a structure-only model and a ten-token rule, and an AUC of 0.75--0.78 when the repository is unseen, with a precision of 2--3\% at the natural base rate. The signal rests on Task-based concurrency, time handling, and resource lifecycles. It fails when the cause of flakiness lies outside the test body, in a helper method, in garbage collector behavior, or in state shared between tests, and when stable tests use time and concurrency in a controlled way.

For practitioners, vocabulary is useful for prioritizing reruns, not for gating pipelines: the top quartile of each repository's ranking contains 23 of the 30 flaky tests. For researchers, the labels, and even the comparison between vocabulary and structure, depend on the rerun round, so rerun studies should resolve tests to their declaring methods, report the sensitivity of their labels, and follow the protocol of Section~\ref{sec:discussion}. As next steps, we plan to make the model follow calls between methods, so that it also reads the code of the helper methods the test uses; to design features that distinguish a fixed wait from time controlled by the test itself; to repeat the study with more projects, chosen at random; \Den{O revisor vai ler isso como ``por que não fizeram agora?''. Os 30 positivos são a principal fraqueza do artigo; se houver tempo antes do ISSTA, ampliar o dataset (mais projetos, rodada em CI/Linux, testes marcados pelos desenvolvedores) é a melhoria de maior retorno.} and to ask practitioners how they would use a test ranking to decide what to rerun.
